\documentclass[10pt,a4paper]{amsart}
\usepackage[margin=19mm]{geometry}
\usepackage{amsmath,amssymb,mathtools}
\usepackage[scr=rsfs,cal=euler]{mathalfa}
\usepackage{xfrac}
\usepackage[unicode,hidelinks,bookmarks=false]{hyperref}
\newcommand{\ie}{i.e.\ }
\newcommand{\cf}{cf.\ }
\newcommand{\resp}{respectively\ }
\newcommand{\Subsection}{\subsection}
\providecommand{\nobreakdash}{\nobreak\hbox{-}}
\setlength{\emergencystretch}{2em}
\multlinegap0pt
\newcommand{\E}{\mathbb E}
\newcommand{\Prob}{\mathbb P}
\newcommand{\R}{\mathbb R}
\newcommand{\Z}{\mathbb Z}
\theoremstyle{plain}
\newtheorem{theorem}{Theorem}[section]
\newtheorem{proposition}[theorem]{Proposition}
\newtheorem{lemma}[theorem]{Lemma}
\newtheorem{corollary}[theorem]{Corollary}
\theoremstyle{remark}
\newtheorem{remark}[theorem]{Remark}
\title[Uniform displacement bounds and Gibbs limits for periodic one-dimensional Riesz gases]{Uniform displacement bounds and Gibbs limits for periodic one-dimensional Riesz gases}
\author{Yan Ru Pei}
\date{Source: arXiv:2609.28503v1 (CC BY 4.0)}
\begin{document}
\maketitle

\noindent\emph{Source and licence.} The delivered work is the article shown in the title above, version arXiv:2609.28503v1 (CC BY 4.0), distributed under the Creative Commons Attribution 4.0 International licence, as recorded in the arXiv licence metadata for this exact version and shown on the abstract page saved with this item. The full licence notice, which carries the licence URL and the address of that abstract page, is the bundled file \texttt{licenses/arxiv-2609.28503-CC-BY-4.0.txt}.
\par\smallskip

\noindent\textit{Editorial scope.} Counted unit: the article's Theorem (source label thm:limit) (source0.tex lines 162--179, source label \texttt{thm:limit}): ``Ordered lattice matching Fix and . The family is relatively compact. Every weak accumulation point along is a simple stationary process of intensity one that admits a representation j j U q j . Here i''. It is delivered together with the article's complete original proof (source0.tex lines 631--720, one proof environment printed later in the article, detached from the statement by the article's own arrangement, 90 lines), copied line for line from the article's own text. Required context retained uncounted: thm:finite (lines 138--150); lem:convexity (lines 325--335). Editorial formatting only: an AMS \texttt{amsart} preamble that restates the article's own macro definitions and its own theorem declarations, and the theorem environments the retained lines use; the article's own class file is neither shipped nor input; the retained environments are renumbered sequentially in order of appearance, whereas the article prints them with its own numbering; the article's equation numbering is not reproduced and every cross-reference inside the retained text resolves to the wrapper's numbering. No citation occurs inside any retained line, so the wrapper carries no bibliography; All retained bodies are exact copies of the bundled \texttt{sources/211593/source0.tex}, so no omission receipt applies. No asset-inclusion command occurs inside any retained span and the package ships no image asset.


\section*{Retained context: thm:finite (source0.tex lines 138--150)}

\begin{theorem}[Uniform particle displacements]\label{thm:finite}
For each $a\in(0,1)$ there is $0<C_a<\infty$ such that, for every
$N\ge2$, $\beta>0$, and $0\le i<N$,
\begin{equation}\label{eq:finite}
 \E Y_i=i,\qquad \E|Y_i-i|^2\le C_a/\beta.
\end{equation}
Set $\sigma=(C_a/\beta)^{1/2}$ and $\kappa=\log(3)/4$.
For every $t\ge0$,
\begin{equation}\label{eq:exp}
 \Prob(|Y_i-i|>t)\le2e^{-\kappa t/\sigma},\qquad
 \E e^{\kappa|Y_i-i|/(2\sigma)}\le3.
\end{equation}
\end{theorem}


\section*{Retained context: lem:convexity (source0.tex lines 325--335)}

\begin{lemma}[Convexity and arc moments]\label{lem:convexity}
The anchored Hamiltonian is strictly convex on $\mathcal D_N$,
and the Palm law is log-concave.  Under either the Palm law or any
of the laws $\E_u$, every cyclic gap has mean one.  If $R$ is the
sum of $k$ consecutive cyclic gaps, $1\le k\le N$, then
\begin{equation}\label{eq:arc-moments}
 \E R=k,\qquad \E R^2\le2k^2,\qquad
 \E R^{2-a}\le2^{(2-a)/2}k^{2-a}.
\end{equation}
The same inequalities hold with $\E_u$ in place of $\E$.
\end{lemma}


\section{The counted unit: Theorem (source label thm:limit) and its complete original proof}

\begin{theorem}[Ordered lattice matching]\label{thm:limit}
Fix $a\in(0,1)$ and $\beta>0$. The family $(\mu_N)$ is relatively
compact. Every weak accumulation point along $N\to\infty$ is a
simple stationary process $\xi$ of intensity one that admits a
representation
\begin{equation}\label{eq:matching}
 \xi=\sum_{j\in\Z}\delta_{j+U+q_j}.
\end{equation}
Here $U$ is uniform on $[0,1)$ and independent of the
$\Z$-stationary sequence $(q_j)$,
$1+q_{j+1}-q_j>0$ almost surely for every $j$, and
\begin{equation}\label{eq:limit-moments}
 \E q_0^2\le C_a/\beta,\qquad
 \Prob(|q_0|>t)\le2e^{-\kappa t/\sigma},\qquad
 \E e^{\kappa|q_0|/(2\sigma)}\le3.
\end{equation}
The correspondence $j+U\mapsto j+U+q_j$ is a stationary matching.
\end{theorem}

\begin{proof}[Proof of Theorem~\ref{thm:limit}]
For the Palm configuration extend $Y_i$ by $Y_{i+N}=Y_i+N$ and
put $p_i=Y_i-i$, so $p_{i+N}=p_i$. Independently choose a uniform
$J\in\{0,\ldots,N-1\}$ and set
\[
 q_j^{(N)}=p_{j+J},\qquad j\in\Z.
\]
This sequence is stationary under integer shifts. With an independent
uniform $U_N\in[0,1)$, define
\[
 y_j^{(N)}=j+U_N+q_j^{(N)}
          =Y_{j+J}+U_N-J.
\]
As a point set this is the periodic Palm configuration translated by
$U_N-J$. The latter variable is uniform modulo $N$ and independent
of the Palm configuration before the label shift. Undoing the Palm
rooting therefore shows that $\sum_j\delta_{y_j^{(N)}}$ has exactly
law $\mu_N$.

Theorem~\ref{thm:finite} gives, for every $N,j$,
\begin{equation}\label{eq:uniform-q}
 \E|q_j^{(N)}|^2\le K,\qquad
 \Prob(|q_j^{(N)}|>t)\le2e^{-\kappa t/\sigma},\qquad
 \E e^{\kappa|q_j^{(N)}|/(2\sigma)}\le3.
\end{equation}
These inequalities follow by averaging over $J$; no log-concavity
of the resulting mixture is needed. Product tightness gives tightness
of $(U_N,q^{(N)})$ in $[0,1]\times\R^{\Z}$. Its weak limits have
uniform $U$, independent of the integer-stationary field $q$.
The moment inequalities pass by lower semicontinuity. For the tail
bound apply the Portmanteau inequality to the open set
$\{|q_j|>t\}$.

The point-process laws themselves are tight: for every compact
interval $B$, stationarity at intensity one gives
$\E\xi_N(B)=|B|$, and Markov's inequality on an increasing sequence
of compact intervals is the usual tightness criterion for locally
finite counting measures. Start with any convergent point-process
subsequence and take a further joint subsequence of
$(U_N,q^{(N)},\xi_N)$. A Skorokhod representation on this Polish
product space makes all these variables converge almost surely.

We identify the limiting point process. For $M>R+2$, a label
$|j|>M$ can enter $[-R,R]$ only if
$|q_j^{(N)}|\ge |j|-R-1$. Hence
\begin{equation}\label{eq:remote-labels}
 \E\sum_{|j|>M}{\bf1}_{\{y_j^{(N)}\in[-R,R]\}}
 \le K\sum_{|j|>M}(|j|-R-1)^{-2}.
\end{equation}
The right side tends to zero uniformly in $N$, and the same bound
holds for $y_j=j+U+q_j$. In particular the limiting sum
$\sum_j\delta_{y_j}$ is locally finite almost surely. For each
compactly supported continuous test function, first restrict the
sum to $|j|\le M$, where coordinate convergence applies, and then
use~\eqref{eq:remote-labels} on the omitted terms. The joint limit
therefore obeys $\xi=\sum_j\delta_{y_j}$ almost surely. This
identification preserves the originally selected marginal limit.

The finite gaps $1+q_{j+1}^{(N)}-q_j^{(N)}$ are positive, so the
limiting gaps are nonnegative. Their marginal law is the common
cyclic gap law of the Palm ensemble. This law is log-concave by
Lemma~\ref{lem:convexity}, has mean one and second moment at most
two. The second-moment bound gives uniform integrability of the
gaps, so every limiting gap has mean one. Weak limits of log-concave
probability measures are log-concave. In one dimension such a law
is either a point mass or has a log-concave density on its affine
support. In the latter case it has no atom at zero; in the former
case its mean forces the mass to lie at one. Consequently every
limiting gap is strictly positive almost surely. A countable
intersection establishes strict ordering for all labels and hence
simplicity.

For a nonnegative compactly supported function $f$, Tonelli's theorem
and stationarity of the field give
\begin{align*}
 \E\sum_{j\in\Z}f(j+U+q_j)
 &=\sum_{j\in\Z}\int_0^1\E f(j+u+q_0)\,du\\
 &=\E\int_\R f(t+q_0)\,dt
 =\int_\R f(t)\,dt.
\end{align*}
Thus the limiting intensity is exactly one.

Finally translate the matching by $t\in\R$. Write
$U+t=m+U'$ with $m\in\Z$ and $U'\in[0,1)$. Reindexing by
$j'=j+m$ sends the translated matching to the representation with
phase $U'$ and field $q'_{j'}=q_{j'-m}$. Conditional on $U$, the
field is independent and integer-stationary, so the law of $q'$ is
the original law and is independent of $U'$. The latter is again
uniform. This proves joint spatial stationarity.
\end{proof}

\end{document}
